\chapter{Hand Tremor}

\section*{Definition}
Hand tremor is an involuntary, rhythmic, oscillatory movement of one or both hands, typically ranging from 4 to 12 Hz in frequency. It may occur at rest (resting tremor), during maintained posture (postural tremor), or with voluntary movement (kinetic/intention tremor), and can significantly impair fine motor tasks such as writing, eating, or dressing.

\section*{Pathophysiology}
Tremor arises from dysfunction within the central nervous system circuits involving the cerebellum, thalamus, basal ganglia, and cortical areas. Essential tremor involves abnormal cerebellar-olivary-thalamic-cortical oscillations. Parkinsonian tremor results from degeneration of dopaminergic neurons in the substantia nigra pars compacta, causing abnormal basal ganglia-thalamic output. Cerebellar tremor results from lesions in the cerebellum or its pathways, disrupting coordination of voluntary movement.

\section*{Causes}
\begin{itemize}
  \item Essential tremor (most common cause of action tremor)
  \item Parkinson disease (resting tremor, pill-rolling, typically asymmetric)
  \item Cerebellar tremor (intention tremor from stroke, MS, tumor, or degeneration)
  \item Drug-induced tremor (beta-agonists, antidepressants, lithium, valproate, amiodarone, corticosteroids)
  \item Metabolic causes (hyperthyroidism, hypoglycemia, hepatic encephalopathy, Wilson disease)
  \item Enhanced physiologic tremor (stress, fatigue, anxiety, caffeine, alcohol withdrawal)
  \item Dystonic tremor (associated with dystonic postures)
  \item Psychogenic tremor (sudden onset, variable, distractible, entrainable)
  \item Orthostatic tremor (leg tremor on standing, may involve hands)
  \item Fragile X-associated tremor/ataxia syndrome (FXTAS)
\end{itemize}

\section*{When to See a Doctor}
See your doctor if:
\begin{itemize}
  \item Tremor interfering with daily activities (writing, eating, drinking, dressing)
  \item Sudden onset or rapid progression of tremor
  \item Tremor accompanied by other neurological symptoms (rigidity, bradykinesia, gait disturbance, speech changes)
  \item Tremor beginning after starting a new medication
  \item Onset before age 40 (may suggest genetic or metabolic etiology)
\end{itemize}

\section*{Related Symptoms}
\begin{itemize}
  \item Rigidity (resistance to passive movement)
  \item Bradykinesia (slowness of movement)
  \item Gait disturbance or imbalance
  \item Changes in handwriting (micrographia in Parkinson, large shaky writing in cerebellar)
  \item Head tremor (yes-yes or no-no oscillation)
  \item Voice tremor (quavering speech)
  \item Fatigue or weakness in affected limb
\end{itemize}

\section*{Self-Care / Home Management}
\begin{itemize}
  \item Ensure a safe environment to prevent injury during episodes; avoid driving or operating machinery if tremor is severe.
  \item Keep a symptom diary documenting triggers, duration, frequency, and associated features.
  \item Practice stress reduction techniques (deep breathing, meditation, adequate sleep) as stress often exacerbates tremors.
  \item Avoid alcohol, caffeine, and recreational drugs which can worsen many tremor conditions.
  \item Seek immediate evaluation for sudden-onset, severe, or progressive neurological symptoms.
\end{itemize}

\section*{How Your Doctor Will Evaluate You}
\begin{itemize}
  \item Detailed neurological history includes onset pattern, progression, triggers, associated symptoms, and functional impact.
  \item Comprehensive neurological examination assesses mental status, cranial nerves, motor/sensory function, reflexes, coordination, and gait.
  \item Neuroimaging (CT, MRI) is often indicated to evaluate for structural, vascular, or demyelinating causes.
  \item Specialized testing (EEG for seizures, nerve conduction studies for neuropathy, lumbar puncture for meningitis) is guided by clinical suspicion.
\end{itemize}

\section*{Treatment Options}
\begin{itemize}
  \item Treatment is directed at the underlying cause identified through diagnostic evaluation.
  \item Symptomatic pharmacotherapy may include beta-blockers for essential tremor, levodopa for Parkinson disease, and antiepileptics for certain tremor types.
  \item Physical, occupational, and speech therapy play important roles in functional recovery and rehabilitation.
  \item Referral to neurology is indicated for unexplained, progressive, or treatment-refractory neurological symptoms.
\end{itemize}

\clearpage
